% GENERATED by scripts/build.py convert ch06 -- do not hand-edit; edit
% chapters/md/06_discussion_and_limitations.md and regenerate.
% md sha256 d8d904483070993c72a51ce16cd334b2a34f4d06d71291dc9e7be9f86a89abad
% canonical manuscript/v13_2026-10-17/06_discussion_and_limitations.md sha256 d8d904483070993c72a51ce16cd334b2a34f4d06d71291dc9e7be9f86a89abad
% md identical to canonical: True
% filter sha256 5f7f0d210d735950ef97c22d7ae51d76d95fce54efc094b16108921e53a47b1c

\section{Discussion and limitations}\label{sec:6}\ThesisCheckNumber{section}{6}

\subsection{Relating the findings to the literature}\label{sec:6.1}\ThesisCheckNumber{subsection}{6.1}

The baseline findings are consistent with the existing literature, but only narrowly. Headey and Palloni \autocite*{HeadeyPalloni2019} find that, under a within-region panel design netting out time-invariant confounding, sanitation improvements predict large reductions in diarrhoea and mortality but are \textbf{not} associated with stunting or wasting, while improved water is statistically insignificant for most outcomes except piped water (reduced stunting). Their null sanitation--stunting finding is a close, specific parallel to this thesis's own null household- and community-sanitation associations with HAZ (Sections 5.1--5.2): both find no detectable sanitation--growth association once a design better suited to netting out confounding is used, each in its own way (region-level panel fixed effects there; region or cluster fixed effects here, with no household-level fixed effect --- household WASH is the exposure, not the FE grouping). This is a narrower, more defensible claim than saying WASH--growth associations ``generally'' attenuate under better controls, which would overstate what any single study, this thesis included, establishes on its own.

The household cluster-fixed-effects associations (Table 5.1) are not significant at the 5\% level, broadly consistent with the strongest randomised evidence available: no WASH-only arm in the three trials showed a statistically detectable growth effect (Section 2.3), while the nutrition arm (no WASH component) showed a modest positive effect, and adding WASH to nutrition brought no additional benefit over nutrition alone \autocite{Cumming2019}. This should not be overstated as confirmation: this thesis's exposure is a facility-type classification in cross-sectional data, not a randomised intervention with prospective follow-up from near birth, so the two literatures answer related but not identical questions, and similarity of findings does not validate this thesis's coefficients. The one significant household-level result (Section 5.3) is negative and just below threshold --- one cell among twelve in that table, not treated as evidence that improved household water harms growth; no mechanism in Section 2 supports that reading, and this thesis proposes none.

The historical extension's pooled water--HAZ association is positive and significant (Section 5.6), independently reproduced with a second estimation library on a separately rebuilt sample, as were all 24 primary established-date cells (Appendix~B). Computational reproduction is not a multiplicity correction: no formal adjustment has been applied for the number of cells estimated across the baseline and historical analyses, so this thesis has neither quantified nor ruled out the possibility that this result is one significant draw among many tested cells. It would be tempting to read the coefficient as evidence that introducing a time dimension recovers an association that the single-round baseline design cannot detect. Section 6.2 explains why this thesis does not draw that conclusion.

\subsection{Why the baseline and historical results cannot be compared as a test of exposure timing}\label{sec:6.2}\ThesisCheckNumber{subsection}{6.2}

This point has its own subsection because it is the most tempting overinterpretation that the structure of this thesis invites. The research question of the historical extension (Sections 1.2 and 1.3) is framed deliberately to avoid it. The extension asks whether modelled early-life local WASH coverage is associated with outcomes \emph{within its own sample}, not whether assigning a historical rather than a current exposure changes an otherwise identical estimate.

The baseline and historical samples differ in at least four respects at the same time, not in one:

\begin{enumerate}
\def\labelenumi{\arabic{enumi}.}
\tightlist
\item
  \textbf{Different children and survey rounds.} The baseline uses 2022--2025 DHS rounds from four countries. The historical extension uses 2014--2018 rounds from three; Ethiopia's 2016 round enters only as a provisional sensitivity analysis outside the established-date pooled sample (Section 1.4). No child appears in both samples.
\item
  \textbf{Different exposure definitions and construction.} The baseline household exposure is the survey response of each child's own household at the survey date. The community exposure is an unweighted leave-one-out proportion among the other sampled, classifiable HR household survey responses in the same cluster, not a census of every household in the cluster (Section 3.3). The historical exposure is a modelled, area-averaged, annual gridded coverage estimate from an external source (IHME), matched to a specific postnatal window (Sections 3.5 to 3.7). These are not the same quantity measured at two different times; they are two different kinds of measurement.
\item
  \textbf{Different eligibility and completion rules.} The historical sample also requires complete spatial coverage and a complete exposure window verified against the interview date (Section 3.7). These restrictions have no counterpart in the baseline sample.
\item
  \textbf{Different control sets and identifying variation.} The main baseline household specification uses cluster fixed effects with household- and individual-level socioeconomic controls, including maternal age and education, literacy, household size and wealth quintile (Section 3.4). In the baseline, region fixed effects appear in the household region-fixed-effects comparison (Section 4.1), in the separate community specification (Section 4.3) and in the joint household--community model (Section 4.4); they are never stacked with the household model's cluster fixed effects. The historical specification uses cluster fixed effects with country-by-birth-year terms and excludes wealth from the main specification, for a reason specific to the historical design (Sections 3.8 and 4.6). The historical cluster-fixed-effects model is identified by genuine within-cluster variation across birth cohorts (Section 4.6). The baseline community model, by contrast, is deliberately \emph{not} estimated with cluster fixed effects, for a reason specific to its leave-one-out construction (Section 4.3) that has no bearing on the historical design.
\end{enumerate}

Because all four differ at once, a contrast between the baseline's largely null associations and the historical extension's selected pooled water--HAZ headline cannot be cleanly attributed to ``the effect of assigning historical instead of current exposure'' --- that headline is one of six nominally significant cells among the 24 primary cells, not the only significant historical result (Section 5.6). The contrast could instead, or in addition, reflect different confounding structures across the two country sets, different measurement-error properties of a modelled surface versus a survey response, the different control sets, or chance given the number of cells estimated (the multiplicity point of Section 5.6 applies here too). No formal test of the baseline-versus-historical difference was carried out, and a significant coefficient beside non-significant ones does not by itself show the two differ significantly. This thesis does not adjudicate between these explanations: its design, two separately constructed samples rather than one shared sample in which only exposure timing varies, cannot do so.

\textbf{What the within-historical-sample window comparison does and does not hold constant.} Section 5.7 compares four exposure windows (birth year, prenatal, post-12-month and post-24-month) within the \emph{same} historical country set and estimation approach. The comparison holds constant the survey rounds, the exposure source and spatial-extraction method, the control set and the estimator. In terms of specification, it changes only which months of the IHME surface are averaged into each child's exposure value. It does \textbf{not} hold the estimation sample constant across the four windows: each window has its own eligible sample, which differs according to which children have complete coverage for that specific window (Section 5.7). Table 5.7 shows how much this can matter. With the window held at post-12-month, the estimate moves from 1.03 on its own eligible sample to 1.64 on the smaller sample eligible under all four windows, a change in sample composition rather than in the window. Differences between the own-sample window estimates therefore combine window changes and sample-composition changes. The window comparison also says nothing about the baseline-versus-historical contrast discussed above, and this paragraph keeps the two separate.

\subsection{Limitations: unresolved documentation}\label{sec:6.3}\ThesisCheckNumber{subsection}{6.3}

Two issues in the historical extension remain genuinely unresolved documentation gaps; they are neither design choices nor inherent limitations. First, no DHS or IPUMS source could be located that documents a measurement-subsample or anthropometry-specific sampling weight, despite repeated attempts with the DHS Biomarker Manual and alternative official sources (Appendix~A records the exact sources tried). This thesis's working inference is that within-survey PERWEIGHT normalisation removes a constant measurement-subsample selection factor, under an assumption of random household-level selection. The inference is stated as such, not as a documented fact, and it cannot correct for differential non-response among children who were eligible but not measured. Second, the gap between Nigeria 2018's official count of measured children (12,806) and the number of children with non-missing anthropometry in the IPUMS birth-history extract (11,704) remains unexplained (Section 3.9).

Both points bear on which children the historical estimates describe, so this thesis reports its historical coefficients as \textbf{weighted conditional associations among eligible measured children}, not population-representative estimates. Reporting the unweighted estimate alongside (Section 5.6) does not resolve this, since the selection question is about who was measured at all, not how measured children are weighted relative to one another; if measurement or non-response related to both local WASH coverage and growth, the estimated associations could differ from the full eligible population, and the data here cannot show whether, or in which direction, they do.

\subsection{Limitations: authorised design choices, disclosed and tested}\label{sec:6.4}\ThesisCheckNumber{subsection}{6.4}

Three features of the design are deliberate, disclosed choices, not errors.

First, the equal-country-total pooled weighting (Section 4.5) gives each established-date country equal total observation weight in a single pooled common-slope regression, not a population-size-weighted estimate, and not the same as averaging each country's own separately estimated coefficient: statistical influence on the pooled slope also depends on residualized exposure variation after controls, so equal weight totals do not imply equal influence. This choice's consequence is disclosed alongside the main estimate throughout: the pooled water--HAZ coefficient moves from 1.03 (weighted) to 0.41 (unweighted, same observations); neither is population-representative (Section 3.8).

Second, the working 12-month postnatal exposure window was chosen during the design phase, before any model was estimated, for a documented reason: it is an early-life window whose completion can be checked against the interview date. It is nowhere described in this thesis as ``preregistered''. The literature basis for this specific window length is modest. Cumming and Cairncross \autocite*{CummingCairncross2016} support early-life exposure in general and, if anything, lean toward a 24-month window, which is why all four candidate windows (Section 5.7) are reported together.

Third, wealth quintile is excluded from the historical extension's main specification because it may be a post-exposure mediator of historical WASH conditions and because it is correlated with the WASH exposure itself (Section 3.8). This is a substantive identification concern. It is not a claim that cluster fixed effects remove wealth mechanically, a claim that the baseline model shows to be false by including and estimating wealth directly (Section 4.2). Section 3.8 also states the competing concern: current wealth may proxy earlier socioeconomic circumstances that confound the association, in which case excluding it leaves that confounding unaddressed. Adding wealth leaves the pooled water--HAZ coefficient essentially unchanged (1.01 vs.~1.03; Section 5.8), but this sensitivity does not adjudicate between the two readings.

\subsection{Limitations: inherent, requiring disclosure rather than resolution}\label{sec:6.5}\ThesisCheckNumber{subsection}{6.5}

DHS/IHME input overlap affects two of the three established-date historical countries, Ghana and Kenya, whose own DHS surveys feed into the IHME gridded estimates used as the exposure. The exposure measure is therefore not fully independent of the survey data used in the analysis. This thesis reports the overlap-exclusion sensitivity alongside the main estimate rather than treating the overlap as removable: the pooled water--HAZ estimate excluding Ghana is 0.74 (Section 5.8). Because that estimate pools Kenya and Nigeria, Kenya's overlap remains in it, so the sensitivity addresses only part of the overlap.

Current residence, as recorded at the survey date, is used as a proxy for residence during each child's early-life window. Children whose survey residence differs from their early-life residence may be assigned an exposure that does not represent their early-life location, under both the primary buffer-average method and the containing-pixel sensitivity. This is a form of exposure measurement error whose extent in these samples is unknown.

DHS cluster coordinates are randomly displaced for confidentiality (Section 3.4), and this limitation holds whatever extraction method is used: the primary buffer average and the containing-pixel sensitivity (Section 5.8) both read pixels around the displaced coordinates, through different mechanisms, and neither recovers the true location or corrects for the undisclosed displacement DHS actually applied (Sections 3.5--3.6).

The IHME exposure surface is itself a \emph{modelled} estimate with its own uncertainty, not a direct observation; its supplied lower and upper bounds are not used as regression inputs, so the standard errors of the historical coefficients do not reflect uncertainty in the exposure surface itself. The surface's annual resolution also means the value used for a window's final touched year can include months after the child's own interview date even once the postnatal completion rule (Section 3.7) is satisfied; the project's construction flags every such case but neither excludes nor adjusts for it (Section 3.7).

The residual within-cluster exposure variation that actually identifies the historical coefficients is small: 0.20 to 0.88 percentage points, depending on country and product (Section 4.6), and 0.43 percentage points for the pooled water--HAZ model (Section 5.6). This is modest relative to the 10-point scale used for reporting, so a coefficient per 10 percentage points is a reporting convention, not a description of a change typically observed within a cluster (Section 5.6). Finally, Ethiopia's interview-date uncertainty (Section 5.9; Appendix~C) is treated as an inherent limitation of that round's recoded date fields, not as a design choice made in this thesis. Its estimates remain provisional and outside the established-date pooled family.

\subsection{What this thesis does not claim}\label{sec:6.6}\ThesisCheckNumber{subsection}{6.6}

This thesis does not claim that its findings establish a causal effect of WASH facilities on child growth, in either direction. Both analyses are cross-sectional and observational, and fixed effects remove only what is shared within a region or cluster; confounding that varies within clusters may remain (Section 4.2). The thesis does not claim that its null or weak baseline associations are evidence that WASH investment is an ineffective use of development resources: non-significance here means that the evidence in this sample and design is inconclusive, not that no association exists. It does not generalise findings from Ethiopia, Ghana, Kenya and Nigeria, four countries selected for data availability and recency (Section 1.4), to Sub-Saharan Africa as a region. It does not treat the historical extension's computational verification (Appendix~B) as evidence of improved causal identification. And, as Section 6.2 sets out, it does not treat the difference between the baseline's largely null results and the historical extension's selected pooled headline result as evidence that exposure timing explains that difference, because the two samples differ in several other respects that the design cannot separate from timing. With these boundaries stated, Section 7 draws together what the evidence of this thesis does, modestly, support.
